% 天体物理学（综述）
% license CCBYSA3
% type Wiki

本文根据 CC-BY-SA 协议转载翻译自维基百科\href{https://en.wikipedia.org/wiki/Astrophysics?utm_source=chatgpt.com}{相关文章}。


天体物理学是一门利用物理学与化学的方法和原理来研究天体及其现象的科学。\(^\text{[2]}\)作为该学科的奠基人之一，詹姆斯·基勒（James Keeler）曾指出，天体物理学“旨在探明天体的本质，而非它们在空间中的位置或运动——研究它们是什么，而非它们在哪里”，\(^\text{[3]}\)这一区别正体现于天体力学的研究范畴中。

天体物理学的研究对象包括太阳（太阳物理学）、其他恒星、星系、系外行星、星际介质以及宇宙微波背景辐射。\(^\text{[4][5]}\)科研人员在整个电磁波谱范围内观测这些天体的辐射，分析的物理特性包括光度、密度、温度及化学成分等。由于天体物理学的研究范围极其广泛，天体物理学家在研究中会综合运用多个物理学分支的概念与方法，包括经典力学、电磁学、统计力学、热力学、量子力学、相对论、核物理与粒子物理，以及原子与分子物理等。

在实践中，现代天文学研究通常涉及大量理论与观测物理学工作。天体物理学的主要研究领域包括：暗物质、暗能量、黑洞及其他天体的性质，以及宇宙的起源与最终命运。\(^\text{[4]}\)
理论天体物理学还研究多个主题，例如：太阳系的形成与演化、恒星动力学与演化、星系的形成与演化、磁流体力学、宇宙大尺度物质结构的形成、宇宙射线的起源、广义相对论与狭义相对论、量子宇宙学与物理宇宙学（即对宇宙最大尺度结构的物理研究），其中也包括弦论宇宙学与天体粒子物理等内容。
\subsection{历史}
天文学是一门古老的科学，长期以来与地球物理学相分离。在亚里士多德的宇宙观中，天体被认为是永恒不变的球体，其唯一的运动是匀速圆周运动；而地上世界则是生长与腐朽的领域，其中的自然运动沿直线进行，并在物体到达其目标时终止。因此，人们认为天界由一种与地球物质根本不同的物质构成——柏拉图认为那是火（Fire），而亚里士多德则称之为以太（Aether）。\(^\text{[6][7]}\)到了十七世纪，伽利略（Galileo）、笛卡尔（Descartes）和牛顿（Newton）等自然哲学家开始主张，天体与地上物体由相似的物质组成，并遵循相同的自然规律。\(^\text{[8][9][10][11]}\)然而，他们所面临的困难在于：当时尚未发明出能够证明这些观点的科学工具。\(^\text{[12]}\)

在十九世纪的大部分时间里，天文学研究的重点仍是例行性的观测工作，如测量天体的位置与计算其运动。\(^\text{[13][14]}\)一种新的天文学——很快被称为天体物理学（Astrophysics）——开始兴起。威廉·海德·沃拉斯顿（William Hyde Wollaston）与约瑟夫·冯·夫琅和费（Joseph von Fraunhofer）分别独立发现，当太阳光被分解时，其光谱中出现了大量暗线（即光强减弱或缺失的区域）。\(^\text{[15]}\)到 1860 年，物理学家古斯塔夫·基尔霍夫（Gustav Kirchhoff）与化学家罗伯特·本生（Robert Bunsen）证明，这些太阳光谱中的暗线与已知气体光谱中的亮线相对应，而每条特定的谱线都对应一种独特的化学元素。\(^\text{[16]}\)基尔霍夫推断，这些暗线是由太阳大气中化学元素吸收光线所造成的。\(^\text{[17]}\)由此，人们首次证明：太阳与恒星中存在的化学元素，也同样存在于地球上。

在进一步研究太阳与恒星光谱的人物中，诺曼·洛基尔（Norman Lockyer）尤为突出。1868 年，他在太阳光谱中发现了辐射线与暗线。与化学家爱德华·弗兰克兰（Edward Frankland）合作研究不同温度与压力下元素的光谱时，他注意到太阳光谱中一条黄色谱线无法对应任何已知元素。于是他提出，这条谱线代表一种新元素，并将其命名为“氦”（Helium），来源于希腊语“Helios”，意为“太阳之神”。\(^\text{[18][19]}\)

1885 年，爱德华·C·皮克林（Edward C. Pickering）在哈佛学院天文台发起了一项雄心勃勃的恒星光谱分类计划。他组织了一支女性计算员团队，其中包括威廉米娜·弗莱明（Williamina Fleming）、安东尼娅·莫里（Antonia Maury）与安妮·跳跃·坎农（Annie Jump Cannon）。她们根据照相底片上记录的光谱对恒星进行了系统分类。到 1890 年，团队已完成超过一万颗恒星的目录，并将它们划分为十三种光谱类型。根据皮克林的设想，到 1924 年，坎农将该目录扩展为九卷本、超过二十五万颗恒星，并发展出“哈佛分类系统”（Harvard Classification Scheme），该系统于 1922 年被国际天文学界正式采纳为全球标准。\(^\text{[20]}\)

1895 年，乔治·埃勒里·黑尔（George Ellery Hale）与詹姆斯·E·基勒（James E. Keeler）联合来自欧洲与美国的十位副主编，共同创办了《天体物理学期刊》（The Astrophysical Journal: An International Review of Spectroscopy and Astronomical Physics）。\(^\text{[21][22]}\)该期刊旨在填补天文学与物理学期刊之间的学术空白，为发表以下研究提供平台：关于光谱仪在天文学中的应用；与天体物理密切相关的实验室研究，包括金属与气体光谱的波长测定、辐射与吸收实验；关于太阳、月球、行星、彗星、流星与星云的理论研究；以及望远镜与实验室仪器的改进与设计。\(^\text{[21]}\)

约在 1920 年，随着赫罗图（Hertzsprung–Russell Diagram）的发现——这一图表至今仍是恒星分类与演化研究的基础——阿瑟·爱丁顿（Arthur Eddington）在论文《恒星的内部结构》（The Internal Constitution of the Stars）中预见了恒星内部核聚变过程的存在与机制。\(^\text{[23][24]}\)当时，恒星能量的来源仍是一个完全的谜团。爱丁顿正确地推测，这一能量来源于氢核聚变为氦的过程，并根据爱因斯坦方程 $E = mc^2$ 解释了能量的释放。这一推断极具前瞻性，因为当时科学界尚未发现核聚变、热核能，甚至尚未明确恒星主要由氢构成（参见“金属丰度”metallicity）。\(^\text{[25]}\)

1925 年，塞西莉亚·海伦娜·佩恩（Cecilia Helena Payne，后为 Cecilia Payne-Gaposchkin）在拉德克利夫学院（Radcliffe College）完成了一篇具有深远影响的博士论文。她将萨哈电离理论（Saha’s Ionization Theory）应用于恒星大气，首次将恒星光谱类型与其温度建立了系统联系。\(^\text{[26]}\)更重要的是，她发现恒星的主要成分是氢与氦，而非与地球相似的化学组成。尽管爱丁顿对此表示支持，但这一结论过于出人意料，以至于她的论文审阅者（包括亨利·诺里斯·罗素 Henry Norris Russell）劝说她在发表前修改了结论。然而，后续研究最终证实了她的发现。\(^\text{[27][28]}\)

到 20 世纪末，对天体光谱的研究已扩展至覆盖从无线电波到可见光、X 射线与伽马射线的广泛波段。\(^\text{[29]}\)进入 21 世纪，观测范围进一步拓展，纳入了基于引力波探测的新型天文观测手段。
\subsection{观测天体物理学}
\begin{figure}[ht]
\centering
\includegraphics[width=6cm]{./figures/fa6be55c6141f1c4.png}
\caption{} \label{fig_AsYs_1}
\end{figure}
观测天文学是天文学的一个分支，主要致力于观测数据的记录与解释；与之相对的是理论天体物理学（Theoretical Astrophysics），后者主要研究物理模型可被测量的推论。  
观测天文学的核心实践是使用望远镜及其他天文仪器来观测天体。

大多数天体物理观测都是通过电磁波谱进行的。
\begin{itemize}
\item 射电天文学研究波长大于数毫米的辐射。典型的研究对象包括：由星际气体与尘埃云等低温天体所发射的射电波；宇宙微波背景辐射——即大爆炸（Big Bang）红移后的光；以及脉冲星（Pulsar），其最初便是在微波频段被探测到的。研究此类射电波需要使用口径巨大的射电望远镜。
\item 红外天文学研究波长过长而无法被肉眼看到、但又短于射电波的辐射。  
红外观测通常使用与可见光望远镜类似的仪器进行。比恒星温度更低的天体（例如行星）通常通过红外波段进行研究。
\item 光学天文学是最早发展的天文学类型。常用仪器包括装配有电荷耦合器件（CCD）或光谱仪（Spectroscope）的望远镜。地球大气会在一定程度上干扰光学观测，因此人们使用自适应光学系统（Adaptive Optics）与空间望远镜（Space Telescopes）来获得尽可能高质量的图像。在这一波段范围内，恒星极为明亮，同时可通过化学光谱分析恒星、星系与星云的化学组成。
\item 紫外线、X 射线与伽马射线天文学研究极高能的天体物理过程，例如双星脉冲星系统（Binary Pulsars）、黑洞（Black Holes）、磁星（Magnetars）等。这类辐射无法有效穿透地球大气，因此科学家采用两种方法观测该波段的电磁辐射：空间望远镜与地面成像大气契伦科夫望远镜（Imaging Air Cherenkov Telescopes, IACT）。前者的代表性天文台包括 RXTE、钱德拉 X 射线天文台（Chandra X-ray Observatory）与康普顿伽马射线天文台（Compton Gamma Ray Observatory）；后者的代表性设施包括高能立体系统（High Energy Stereoscopic System, H.E.S.S.）与 MAGIC 望远镜。
\end{itemize}
除电磁辐射之外，从地球上能够观测到的、起源于遥远宇宙的信号极为有限。目前已建成少数引力波观测台，但引力波（Gravitational Waves）的探测极其困难。此外，人们还建造了中微子观测台（Neutrino Observatories），主要用于研究太阳。由极高能粒子组成的宇宙射线（Cosmic Rays）也可以通过其撞击地球大气层时产生的效应进行观测。


观测的时间尺度也可能存在显著差异。  
大多数光学观测的时间范围为数分钟至数小时，因此变化速度更快的天体现象往往难以直接观测。  
然而，一些天体具有跨越数百年甚至数千年的历史观测记录。  
另一方面，射电观测既可以分析毫秒量级的事件（如毫秒脉冲星，\textit{Millisecond Pulsars}），也可以结合多年数据进行分析（例如脉冲星减速研究）。  
这些不同时间尺度的观测所得信息差异极大。

\paragraph{}
太阳的研究在观测天体物理学中占据特殊地位。  
由于其他恒星距离极为遥远，太阳是唯一能够以无与伦比的细节进行直接观测的恒星。  
理解太阳的结构与活动机制，有助于我们理解其他恒星的物理特性与演化规律。

\paragraph{}
关于恒星如何变化，即恒星演化（\textit{Stellar Evolution}）的研究，通常通过将不同类型的恒星放置在赫罗图（\textit{Hertzsprung–Russell Diagram}）上加以建模。  
该图可视为描述恒星从诞生到毁灭整个生命周期状态的示意图。
